\subsection{Dimension and Coordinates}

Throughout this subsection, let $\mathcal{V}$ be a vector space over $\mathbb{F}$.

\begin{theorem} \label{thm:basis-length-well-defined} \hlt{Basis Length Is Well-Defined}\\
If $\mathcal{V}$ is finite-dimensional and $\beta_1,\beta_2$ are bases of $\mathcal{V}$ of lengths $m,n$, respectively, then $m=n$.
\begin{proof}
The inequality $m\leq n$ holds by Theorem \ref{thm:finitedimlenofvec}, with $\beta_1$ independent and $\beta_2$ spanning.\\
After interchanging the bases, $n\leq m$ holds by the same theorem. Hence $m=n$.\\
\hlt{Structural Commentary}: Using both basis roles proves the opposite length inequalities.\\
\hlt{Key Insight}: Symmetric comparison makes basis length intrinsic.
\end{proof}
\end{theorem}

\begin{definition} \label{def:dimension} \hlt{Dimension of a Finite-Dimensional Space}\\
Assume $\mathcal{V}$ is finite-dimensional. Choose a basis $\beta=(\mathbf{v}_1,\ldots,\mathbf{v}_n)$ by Corollary \ref{cor:finite-dimensional-spaces-and-subspaces-have-bases}, and define $\dim(\mathcal{V})\defeq n$.\\
This value is independent of $\beta$ by Theorem \ref{thm:basis-length-well-defined}, and $\mathcal{V}$ is called \hlt{$n$-dimensional}. Since $()$ is a basis of $\{\mathbf{0}\}$, $\dim(\{\mathbf{0}\})=0$.
\end{definition}

\begin{proposition} \label{prop:dimension-bounds-and-basis-criteria} \hlt{Subspace Dimension Criterion}\\
Let $\mathcal{V}$ be finite-dimensional and let $\mathcal{U}\leq\mathcal{V}$. Then:
\begin{enumerate}[label=(\roman*)]
\setlength{\itemsep}{0pt}
\item $\dim(\mathcal{U})\leq\dim(\mathcal{V})$;
\item $\dim(\mathcal{U})=\dim(\mathcal{V})$ iff $\mathcal{U}=\mathcal{V}$.
\end{enumerate}
\begin{proof}
Choose a basis $\delta=(\mathbf{u}_1,\ldots,\mathbf{u}_m)$ of $\mathcal{U}$ and a basis $\beta=(\mathbf{v}_1,\ldots,\mathbf{v}_n)$ of $\mathcal{V}$ by Corollary \ref{cor:finite-dimensional-spaces-and-subspaces-have-bases}.
\begin{enumerate}[label=(\roman*)]
\setlength{\itemsep}{0pt}
\item[$(i)$] The list $\delta$ is independent in $\mathcal{V}$, while $\beta$ spans $\mathcal{V}$. Hence $m\leq n$ by Theorem \ref{thm:finitedimlenofvec}, so $\dim(\mathcal{U})\leq\dim(\mathcal{V})$.
\item[$(ii)$] ($\Rightarrow$) Suppose $\dim(\mathcal{U})=\dim(\mathcal{V})$. Then $m=n$. Extend $\delta$ to a basis of $\mathcal{V}$ by Theorem \ref{thm:extending-independent-lists}.\\
The extension adds no entry because its length is $n$ by Theorem \ref{thm:basis-length-well-defined}. Thus $\mathcal{U}=\operatorname{span}(\delta)=\mathcal{V}$.
\item[$(ii)$] ($\Leftarrow$) If $\mathcal{U}=\mathcal{V}$, then $\delta$ is a basis of $\mathcal{V}$, so $m=n$ by Theorem \ref{thm:basis-length-well-defined}. Hence $\dim(\mathcal{U})=\dim(\mathcal{V})$.
\end{enumerate}
\hlt{Structural Commentary}: Length comparison proves the bound; basis extension tests equality against a new ambient direction.\\
\hlt{Key Insight}: A subspace of full ambient dimension cannot be proper.
\end{proof}
\end{proposition}

\begin{proposition} \label{prop:full-dimension-basis-criteria} \hlt{Full-Dimension Basis Criteria}\\
Let $\mathcal{V}$ be finite-dimensional and set $n\defeq\dim(\mathcal{V})$. Then:
\begin{enumerate}[label=(\roman*)]
\setlength{\itemsep}{0pt}
\item if $\gamma$ is a spanning list in $\mathcal{V}$ of length $n$, then $\gamma$ is a basis;
\item if $\alpha$ is a linearly independent list in $\mathcal{V}$ of length $n$, then $\alpha$ is a basis.
\end{enumerate}
\begin{proof}
\begin{enumerate}[label=(\roman*)]
\setlength{\itemsep}{0pt}
\item[$(i)$] A basis sublist $\beta$ of $\gamma$ exists by Theorem \ref{thm:reducing-spanning-lists}. Its length is $n$ by Definition \ref{def:dimension} and Theorem \ref{thm:basis-length-well-defined}.\\
Since $\gamma$ also has length $n$, the sublist $\beta$ equals $\gamma$. Hence $\gamma$ is a basis.
\item[$(ii)$] Extend $\alpha$ to a basis $\beta$ of $\mathcal{V}$ by Theorem \ref{thm:extending-independent-lists}. The length of $\beta$ is $n$ by Definition \ref{def:dimension}.\\
Hence the extension adds no entry to the length-$n$ list $\alpha$. Hence $\alpha=\beta$ is a basis.
\end{enumerate}
\hlt{Structural Commentary}: At full length, basis reduction cannot delete and basis extension cannot append.\\
\hlt{Key Insight}: At the ambient dimension, either half of the basis condition forces the other.
\end{proof}
\end{proposition}

\begin{theorem} \label{thm:arbitrary-bases-same-cardinality} \hlt{Dimension as Basis Cardinality}\\
Assume the Choice Principle \ref{princ:axiom-of-choice-analysis}. If $\mathcal{B}$ and $\mathcal{C}$ are Hamel bases of $\mathcal{V}$, then $\abs{\mathcal{B}}=\abs{\mathcal{C}}$.
\begin{proof}
\begin{enumerate}[label=(\roman*)]
\setlength{\itemsep}{0pt}
\item \hlt{Finite case}: Assume $\abs{\mathcal{B}}=n\in\N_0$, and order $\mathcal{B}$ as a spanning list $\beta$. Then $\mathcal{V}$ is finite-dimensional.\\
The set $\mathcal{C}$ cannot contain $n+1$ distinct vectors by Theorem \ref{thm:finitedimlenofvec}, because they would form an independent list longer than $\beta$.\\
Thus $\mathcal{C}$ is finite; otherwise finite recursion would select such vectors. If $\gamma$ orders $\mathcal{C}$ and has length $k$, then $k=n$ by Theorem \ref{thm:basis-length-well-defined}, and $\abs{\mathcal{C}}=k=n$ by Proposition \ref{prop:finite-cardinality-and-aleph-zero-agreement}(i). Interchanging the bases handles finite $\mathcal{C}$.
\item \hlt{Support incidence}: Assume $\mathcal{B}$ and $\mathcal{C}$ are infinite. For $\mathbf{c}\in\mathcal{C}$, let $a_{\mathbf{c}}:\mathcal{B}\to\mathbb{F}$ be its unique finite $\mathcal{B}$-expansion by Proposition \ref{prop:hamel-coordinate-representation}, and set $S_{\mathbf{c}}\defeq\operatorname{supp}_{\mathcal{B}}(a_{\mathbf{c}})$. Fix $\mathbf{b}\in\mathcal{B}$ and write $\mathbf{b}=\sum\limits_{j=1}^{m}d_j\mathbf{c}_j$ with $\mathbf{c}_j\in\mathcal{C}$.\\
If $\mathbf{b}\notin\bigcup\limits_{\mathbf{c}\in\mathcal{C}}S_{\mathbf{c}}$, then substituting the $\mathcal{B}$-expansions makes the coefficient at $\mathbf{b}$ equal to $\sum\limits_{j=1}^{m}d_j a_{\mathbf{c}_j}(\mathbf{b})=0$, whereas the standard expansion of $\mathbf{b}$ has coefficient $1$. This discrepancy is impossible by uniqueness in Proposition \ref{prop:hamel-coordinate-representation}.\\
Hence $(\forall\,\mathbf{b}\in\mathcal{B})(\exists\,\mathbf{c}\in\mathcal{C})\,\mathbf{b}\in S_{\mathbf{c}}$. Set $D\defeq\{(\mathbf{c},\mathbf{b})\in\mathcal{C}\times\mathcal{B}:\mathbf{b}\in S_{\mathbf{c}}\}$.
\item \hlt{Cardinal comparison}: Choose $(\mathbf{c}_{\mathbf{b}})_{\mathbf{b}\in\mathcal{B}}$ such that $(\forall\,\mathbf{b}\in\mathcal{B})\,[\mathbf{c}_{\mathbf{b}}\in\mathcal{C}\land\mathbf{b}\in S_{\mathbf{c}_{\mathbf{b}}}]$ by Principle \ref{princ:axiom-of-choice-analysis}.\\
Then $\mathbf{b}\mapsto(\mathbf{c}_{\mathbf{b}},\mathbf{b})$ injects $\mathcal{B}$ into $D$. Fix a well-order of $\mathcal{B}$ by Principle \ref{princ:ordinal-foundations}(iv), let $p_{\mathbf{c}}(\mathbf{b})\in\N$ be position of $\mathbf{b}$ in the induced ordering of the finite set $S_{\mathbf{c}}$. The map $(\mathbf{c},\mathbf{b})\mapsto(\mathbf{c},p_{\mathbf{c}}(\mathbf{b}))$ injects $D$ into $\mathcal{C}\times\N$. Hence $\abs{\mathcal{B}}\leq\abs{D}\leq\abs{\mathcal{C}\times\N}$ by Proposition \ref{prop:cardinal-operations-on-representatives}(i).\\
By Proposition \ref{prop:cardinal-operations-on-representatives}(iv) and Proposition \ref{prop:finite-cardinality-and-aleph-zero-agreement}(iii), $\abs{\mathcal{C}\times\N}=\abs{\mathcal{C}}\aleph_0$.\\
Since $\mathcal{C}$ is infinite, $\aleph_0\leq\abs{\mathcal{C}}$ by Proposition \ref{prop:finite-cardinality-and-aleph-zero-agreement}(iv), so $\abs{\mathcal{C}}\aleph_0=\abs{\mathcal{C}}$ by Theorem \ref{thm:infinite-cardinal-absorption-reference}(ii). Therefore $\abs{\mathcal{B}}\leq\abs{\mathcal{C}}$.
\item \hlt{Symmetry}: Interchanging $\mathcal{B}$ and $\mathcal{C}$ proves $\abs{\mathcal{C}}\leq\abs{\mathcal{B}}$. These inequalities correspond to injections in both directions by Proposition \ref{prop:cardinal-operations-on-representatives}(i), so $\mathcal{B}\sim_{\#}\mathcal{C}$ by Theorem \ref{thm:schroeder-bernstein}. Hence $\abs{\mathcal{B}}=\abs{\mathcal{C}}$ by Proposition \ref{prop:cardinality-canonical-and-bijection-invariant}(ii).
\end{enumerate}
\hlt{Structural Commentary}: Finite supports form an incidence relation that choice and well-ordering encode by cardinal injections.\\
\hlt{Key Insight}: Record each support occurrence by one basis vector and one finite position.
\end{proof}
\end{theorem}

\begin{definition} \label{def:hamel-dimension} \hlt{Hamel Dimension}\\
Assume Zorn's lemma and the Choice Principle \ref{princ:axiom-of-choice-analysis}. Choose a Hamel basis $\mathcal{B}$ of $\mathcal{V}$ by Theorem \ref{thm:hamel-basis-extension}, and define $\dim(\mathcal{V})\defeq\abs{\mathcal{B}}$.\\
The value is independent of $\mathcal{B}$ by Theorem \ref{thm:arbitrary-bases-same-cardinality}. If $\mathcal{V}$ has a finite ordered basis of length $n$, its underlying set is a Hamel basis of cardinality $n$ by Proposition \ref{prop:finite-cardinality-and-aleph-zero-agreement}(i); hence this definition extends Definition \ref{def:dimension}.
\end{definition}
